% ============================================================================
%  CLEANED / CORRECTED edition of:
%
%      George Kempf and Linda Ness,
%      "The length of vectors in representation spaces",
%      in: Algebraic Geometry (Copenhagen 1978),
%      Lecture Notes in Mathematics 732, Springer-Verlag, pp. 233-243.
%      DOI 10.1007/BFb0066647
%
%  This is a faithful transcription with spelling/transcription errors and a
%  handful of genuine errata of the original typescript corrected.  Every change
%  relative to the verbatim transcription (kempf-ness-verbatim.tex) is flagged
%  inline with a "% CORRECTION" comment.  Mathematical statements were checked
%  against Mumford, Geometric Invariant Theory (the notion of a "properly
%  stable" point requires a FINITE stabilizer), and against the paper's own
%  internal usage.
%
%  Summary of corrections:
%   1. intro: "algebriac" -> "algebraic"
%   2. heading: mis-numbered "\S4. The reductive case" -> "\S3."
%   3. p.238: "by che results" -> "by the results"
%   4. p.239: statement *), "such that that" -> "such that the"
%   5. p.239: defn of properly stable, "G_v is connected" -> "G_v is finite"
%   6. p.240: "criterian" -> "criterion"
%   7. p.242: "the only possiblity" -> "possibility"
%   8. p.242: "consists the elements" -> "consists of the elements"
%   9. ref [1]: "algebriac" -> "algebraic"
%  10. ref [3]: "arithmetiques" -> "arithmetiques" with accent
%  11. ref [6]: "Instability invariant theory ... to appear" -> correct title
%      "Instability in invariant theory" with published coordinates
%  Two scan-restorations (already mathematically forced) are also present:
%   - p.239 c): "...is closed and v != 0" (per the introduction's definition)
%   - p.241: the relevant limit is x -> -infinity (since x = log|t| and t -> 0)
% ============================================================================
\documentclass[11pt]{article}

\usepackage[T1]{fontenc}
\usepackage[utf8]{inputenc}
\usepackage{amsmath}
\usepackage{amssymb}
\usepackage{amsthm}
\usepackage{enumitem}
\usepackage[margin=1in]{geometry}

% --- theorem-like environments with the original (manual) numbering ----------
\theoremstyle{plain}
\newtheorem{thm}{Theorem}
\newtheorem{lem}{Lemma}
\newtheorem{prop}{Proposition}
\newtheorem{cor}{Corollary}

\renewcommand{\qedsymbol}{Q.E.D.}

\newlist{abc}{enumerate}{1}
\setlist[abc]{label=\textup{\alph*)}, leftmargin=*, labelsep=0.6em, itemsep=0.4ex}

\newcommand\blfootnote[1]{%
  \begingroup
  \renewcommand\thefootnote{}\footnote{#1}%
  \addtocounter{footnote}{-1}%
  \endgroup
}

\newcommand{\norm}[1]{\lVert #1 \rVert}
\newcommand{\C}{\mathbb{C}}
\newcommand{\R}{\mathbb{R}}
\newcommand{\Z}{\mathbb{Z}}

\begin{document}

\begin{center}
  {\large The length of vectors in representation spaces}\\[2.5em]
  \begin{tabular}{c@{\hspace{2.5em}}c@{\hspace{2.5em}}c}
    George Kempf\textsuperscript{*} & and & Linda Ness\\[0.3em]
    The Johns Hopkins University &  & University of Washington\\
    and Princeton University &  & and The Institute for\\
     &  & Advanced Study
  \end{tabular}
\end{center}
\blfootnote{*Partly supported by NSF contract \#MPS75-05578.}

\bigskip

Let $V$ be a finite dimensional morphic representation of a connected reductive
% CORRECTION: "algebriac" -> "algebraic"
algebraic group $G$ over the complex numbers.
Given a maximal compact subgroup $K$ of $G$, we will fix a Hermitian norm
$\norm{\ }$ on $V$ so that the action of $K$ on $V$ preserves this norm.

Let $v$ be a vector in $V$. The purpose of this paper is to study how the length
changes as one moves along the orbit $G\cdot v$. Thus we want to examine the
function $p_v(g)=\norm{g\cdot v}^2$ on $G$. If $G_v$ denotes the stabilizer of
$v$ in $G$, the function $p_v$ on $G$ is invariant on the left by $K$ and on the
right by $G_v$. Hence $p_v$ is constant on double $K-G_v$ cosets of the form
$K\cdot g\cdot G_v$.

In this paper we will show that the function $p_v$ has very special properties.
We will first prove

\renewcommand{\thethm}{0.1}
\begin{thm}\leavevmode
\begin{abc}
\item Any critical point of $p_v$ is a point where $p_v$ obtains its minimum
      value.
\end{abc}
If $p_v$ obtains a minimum value,
\begin{abc}[start=2]
\item then the set $m$ where $p_v$ obtains this value consists of a single
      $K-G_v$ coset and is connected, and
\item the second order variation of $p_v$ at a point of $m$ in any direction not
      tangent to $m$ is positive.
\end{abc}
\end{thm}

With this theorem in mind, one may ask, ``when does $p_v$ obtain a minimum
value?'' If $v$ is a stable vector (i.e., the orbit $G\cdot v$ is closed and
$v\neq 0$), then $p_v$ clearly obtains a minimum value. The converse is also true
by

\renewcommand{\thethm}{0.2}
\begin{thm}
The vector $v$ is stable if and only if $p_v$ obtains a minimum value.
\end{thm}

The development of the properties of the function $p_v$ is similar to that of
Mumford's numerical function (See [8] and [6]). Also the functions $p_v$ for
certain unstable vectors $v$ play an interesting role in Borel's treatment of
reduction theory [3]. Furthermore, related material may be found in [1,4,7]. We
hope that our results may be useful for studying moduli via geometric invariant
theory.

In section one, we will deal abstractly with a more general type of function
related to the functions like $p_v$. In the next section, we will apply the
previous ideas to prove Theorem 0.1 when $G$ is an algebraic torus. In the third
section, we will treat a general reductive group. In the last section, we will
prove Theorem 0.2.

\bigskip
\noindent\textbf{\S1.\quad Special functions on affine spaces}
\medskip

Let $A$ be a finite dimensional affine space over the real numbers. A special
function on $A$ is a finite sum $\sum e^{H_i(a)}$ where the $H_i$ are affine
functions on $A$.

We will begin by studying special functions where $A$ is the real line $\R$.
Thus a special function on $\R$ may be written uniquely as
$f(x)=\sum a_i e^{\ell_i\cdot x}$, where the $a_i$'s are positive real numbers
and the $\ell_i$'s are distinct real numbers. By calculus, we may deduce

\renewcommand{\thelem}{1.1}
\begin{lem}\leavevmode
\begin{abc}
\item The second derivative $f''$ is never negative.
\item If $f''$ is zero anywhere, then $f$ is constant.
\item A non-constant special function on $\R$ is a strictly convex Morse
      function.
\end{abc}
\end{lem}

\begin{proof}
For a), note that $f''(x)=\sum (a_i\cdot\ell_i^{2})e^{\ell_i\cdot x}$ is the sum
of non-negative terms and, hence, $f''$ is never negative. For b), if
$f''(x)=0$ for some $x$, then $a_i\cdot\ell_i^{2}=0$ for each $i$. This can only
happen if $\ell_i=0$ as $a_i\neq 0$. Therefore, if $f''(x)=0$ for some $x$, then
$f$ is a constant function. This proves b). Part c) results from a) and b).
\end{proof}

A special function on an affine space is called degenerate if its restriction to
any line is constant. We may generalize part a) of the last lemma in

\renewcommand{\theprop}{1.2}
\begin{prop}
Any non-degenerate special function $f$ on an affine space is a strictly convex
Morse function. Hence, if $f$ has a critical point, it must be the unique point
where $f$ obtains its minimum value.
\end{prop}

\begin{proof}
As $f$ is non-degenerate, its restriction to any line is a strictly convex Morse
function. Therefore, $f$ is strictly convex and has positive second order
variation through its critical points. This proves the first statement. The
second statement follows formally from the first.
\end{proof}

To finish our discussion of special functions, we will show how to reduce the
general case to the non-degenerate case. The reduction will be achieved by means
of

\renewcommand{\thelem}{1.3}
\begin{lem}
Let $f$ be a special function on an affine space $A$. There is a unique quotient
affine space $B$ of $A$ and a non-degenerate special function on $B$ such that
$f=g\circ\pi$ where $\pi:A\to B$ is the quotient mapping.
\end{lem}

\begin{proof}
If $f=\sum e^{H_i(a)}$ where the $H_i$'s are affine functions on $A$, there is a
maximal affine subspace $M_a$ through any point $a$ of $A$ such that each
$H_i|_{M_a}$ is constant. As the $M_a$'s are parallel and have the same
dimension, these subspaces $\{M_a\}$ form a quotient affine space, say $B$.
Clearly, $f$ descends to a special function $g$ on $B$. One may easily check that
$g$ is non-degenerate by using the maximality of the $M_a$. The uniqueness
assertion is obvious.
\end{proof}

We may now generalize the last proposition in

\renewcommand{\thethm}{1.4}
\begin{thm}
Let $f$ be a special function on an affine space $A$. Then,
\begin{abc}
\item $f$ is convex,
\item $f$ has no critical points outside of the set $m$ of points where $f$
      obtains its minimum value,
\item $m$ is an affine subspace of $A$,
\item at any point of $m$, the second order variation of $f$ is positive in any
      direction not tangent to $m$.
\end{abc}
\end{thm}

\begin{proof}
By Lemma 1.3, the stated properties may be easily deduced from the Proposition
1.2.
\end{proof}

\bigskip
\noindent\textbf{\S2.\quad The toroidal case.}
\medskip

In this section, we will assume that the group $G$ of the introduction is an
algebraic torus $T=\{(t_1,\dots,t_n)\mid t_i\in\C-\{0\}\}$. The maximal compact
subgroup $K_T$ consists of the elements of $T$ whose coordinates have absolute
value one.

Consider the morphic representation of $T$ on $V$. Then $V$ may be written
uniquely as the direct sum $V=\bigoplus V_\chi$ where $V_\chi$ is the
$\chi$-eigenspace of $V$ for the characters $\chi$ of $T$. Recall that a
character $\chi:T\to\C-\{0\}$ is a (morphic) homomorphism sending
$(t_1,\dots,t_n)$ to $\prod t_i^{m_i}$ where $(m_1,\dots,m_n)$ is a sequence of
integers.

Our Hermitian norm $\norm{\ }$ on $V$ is invariant under $K_T$ if and only if two
eigenspaces with distinct characters are perpendicular. With this information, we
may easily determine the nature of the functions $p_v(t)=\norm{t\cdot v}^2$ for
any vector $v$ in $V$.

\renewcommand{\thelem}{2.1}
\begin{lem}
For any vector $v$ in $V$, there is a finite set $\Xi(v)$ of elements in $\Z^n$
such that
\begin{abc}
\item $p_v((t_1,\dots,t_n))=\displaystyle\sum_{(m_i)\in\Xi(v)}(\text{some positive real})
      \prod|t_i|^{2m_i}$ \quad and
\item the stabilizer $T_v$ of $v$ in $T$ is given by the system of equations
      \[ \prod t_i^{m_i}=1 \quad\text{for } (m_i)\in\Xi(v). \]
\end{abc}
\end{lem}

\begin{proof}
Let $v=\sum v_\chi$ be the eigendecomposition of $v$ where each $v_\chi$ is
non-zero. Let $\Xi(v)$ be the set of sequences $(m_i)$ corresponding to the
characters which occur in this decomposition. As
$(t_1,\dots,t_n)\cdot v=\sum_{(m_i)\in\Xi(v)}\prod t_i^{m_i}\cdot v_{(m_i)}$,
the part b) is evident. Using the orthogonality mentioned above, we have
\[
p_v(t_1,\dots,t_n)=\norm{(t_1,\dots,t_n)\cdot v}^2
=\sum\Bigl\lVert\textstyle\prod t_i^{m_i}\cdot v_{(m_i)}\Bigr\rVert^2
=\sum\norm{v_{(m_i)}}^2\prod|t_i|^{2m_i}
\]
This proves a).
\end{proof}

To prove the desired facts about the function $p_v$, we may study the induced
function $p_v'$ on the quotient Lie group $K_T\backslash T$. Next we will
introduce coordinates on the quotient $K_T\backslash T$. Define
$x(t_1,\dots,t_n)=(\log|t_1|,\dots,\log|t_n|)=(x_1,\dots,x_n)$ for any
$(t_1,\dots,t_n)$ in $T$. Thus $x$ defines an isomorphism from $K_T\backslash T$
with the vector space $\R^n$. With these coordinates we may reinterpret the last
lemma in

\renewcommand{\thelem}{2.2}
\begin{lem}
For any vector in $V$,
\begin{abc}
\item $p_v'(x_1,\dots,x_n)=\displaystyle\sum_{(m_i)\in\Xi(v)}
      e^{\,\text{real}\,+\,2\sum m_i x_i}$ \quad and
\item the locus $\{\sum m_i x_i=0\text{ for all }(m_i)\text{ in }\Xi(v)\}$ is the
      image of $T_v$ in $K_T\backslash T$
\end{abc}
\end{lem}

\begin{proof}
This is a trivial reformulation of Lemma 2.1.
\end{proof}

We will now apply the theory of section one to the function $p_v'$ on
$K_T\backslash T$. By part a) of Lemma 2.2, we know that $p_v'$ is a special
function on the vector space $K_T\backslash T=\R^n$. By the proof of Lemma 1.3
and the part b) of Lemma 2.2, we know that the function $p_v'$ comes from a
non-degenerate special function on
\[
\R^n/\{\textstyle\sum m_i x_i=0\mid(m_i)\in\Xi(v)\}=K_T\backslash T/T_v.
\]
Therefore, by the results of the last section, we have verified the theorem 0.1
% CORRECTION: "by che results" -> "by the results"
when $G$ is an algebraic torus. We will next show how the general case follows
from the toroidal case.

\bigskip
\noindent\textbf{\S3.\quad The reductive case.}% CORRECTION: original mis-numbers this "\S4"; the introduction calls it the third section
\medskip

Using E.~Cartan's results about the geometry of $K\backslash G$, we will be able
to deduce the reductive case from the toroidal case of Theorem 0.1. First we will
recall some of Cartan's results.

Let $\mathcal{T}$ be the set of maximal algebraic tori $T$ in the reductive group
$G$ such that $K\cap T$ is the maximal compact subgroup $K_T$ of $T$. Thus
$K_T\backslash T$ is a submanifold of $K\backslash G$. We will need to know the
Cartan decomposition,
\begin{abc}[label=\textup{\arabic*)}]
\item $K\backslash G=\bigcup_{T\in\mathcal{T}}K_T\backslash T$

      and its infinitesimal form
\item tangent spaces of $K\backslash G$ at $K=\bigcup_{T\in\mathcal{T}}$ tangent
      space of $K_T\backslash T$ at $K_T$.
\end{abc}
[\,for instance, 5\,].

To prove part a) of theorem 0.1, it will suffice to assume that
\begin{description}[leftmargin=2em]
\item[A)] $e$ is the critical point of $p_v$, as $p_{h\cdot v}(g)=p_v(g\cdot h)$.
\end{description}
Let $g$ be any point of $G$. We want to prove that $p_v(g)\geq p_v(e)$. By 1),
$g=k\cdot t$ where $k\in K$ and $t\in T$ for some $T$ in $\mathcal{T}$. By the
left $K$-invariance of $p_v$, we have $p_v(g)=p_v(t)$. As $e$ must be a critical
point of $p_v|_T$, we have the inequality $p_v(t)\geq p_v(e)$ by the toroidal
case of the theorem. Hence, $p_v(g)\geq p_v(e)$ which proves part a) of Theorem
0.1.

To prove the last two parts of the theorem, it will suffice to treat the
case where
\begin{description}[leftmargin=2em]
\item[B)] $p_v$ obtains its minimum value at $e$ (i.e., $K\cdot G_v\subseteq m$),
          because $p_{h\cdot v}(g)=p_v(g\cdot h)$ and $G_{h\cdot v}=hG_vh^{-1}$.
\end{description}
Let $g$ be an element of $m$. If we write $g=k\cdot t$ as before, we find that
$t$ is a point of $T$ where $p_v|_T$ obtains its minimum value $p_v(e)$. By the
toroidal case, $t\in K_T\cdot T_v$ which is a connected subset containing $e$.
Therefore, $m\subseteq K\cdot(\bigcup_{T\in\mathcal{T}}K_T\cdot T_v)\subseteq
K\cdot G_v$ and $\bigcup_{T\in\mathcal{T}}K_T\cdot T_v$ is connected. Hence,
$m=K\cdot G_v\,(=K\cdot\bigcup_{T\in\mathcal{T}}K_T\cdot T_v)$ is connected
because $K$ is also connected as $G$ is. This proves part b) of the theorem 0.1.

To prove part c) of theorem 0.1, it again will suffice to prove the statement
about directions through the point $e$ of $m$. Let $p_v'$ be the function induced
by $p_v$ on $K\backslash G$. We want to show that
\begin{description}[leftmargin=2em]
\item[*)] If $X$ is a tangent vector to $K\backslash G$ at $K$ such that the
          second order variation of $p_v'$ along $X$ is zero, then $X$ is tangent
          to $K\backslash K\cdot G_v$.% CORRECTION: "such that that" -> "such that the"
\end{description}
By 2), $X$ is tangent to $K_T\backslash T$ for some $T$ in $\mathcal{T}$. By the
toroidal case of the theorem, $X$ must be tangent to
$K_T\backslash K_T\cdot T_v$ which is contained in $K\backslash K\cdot G_v$.
This completes the proof of theorem 0.1.

\bigskip
\noindent\textbf{\S4.\quad Analysis of stability}
\medskip

In this section we will study the relationship between the stability of a vector
$v$ in $V$ and the function $p_v$. We will exclude the trivial case $v=0$.

Recall the definitions of various notions of stability. The vector is called
\begin{abc}
\item unstable if $\overline{G\cdot v}\ni 0$,
\item not stable if $G\cdot v$ is not closed,
\item stable if $G\cdot v$ is closed and $v\neq 0$,
\item properly stable if $v$ is stable and $G_v$ is finite.% CORRECTION: original says "connected"; a properly stable point has a FINITE stabilizer (Mumford, GIT), as the proofs below also require
\end{abc}
Trivially one has the following criterion:% CORRECTION: "criterian" -> "criterion"
\begin{description}[leftmargin=2em]
\item[($\dagger$)] $v$ is unstable $\iff \inf p_v=0$.
\end{description}

The Theorem 0.2 will give us a criterion for $v$ being stable in terms of the
function $p_v$. To prove Theorem 0.2, we have to show that,
\begin{description}[leftmargin=2em]
\item[($*$)] if $v$ is not stable, then $\inf p_v$ is not a value of $p_v$.
\end{description}

\begin{proof}[Proof of Theorem 0.2.]
Regard $G$ as an algebraic group defined over $\R$ such that $K$ is the real
locus of $G$. Then the tori $T$ in $\mathcal{T}$ are exactly the maximal tori of
$G$, which are defined over $\R$.

Assume that the vector $v$ is not stable. Then we may find a parabolic subgroup
$P$ of $G$ such that each maximal torus $T$ of $P$ contains a one-parameter
algebraic subgroup $\lambda_T:\C-\{0\}\to T$ such that
$\lim_{t\to 0}\lambda(t)\cdot v$ exists in $V$ and is a point outside the orbit
$G\cdot v$. See [1] or [6].

Let $\overline{P}$ be the parabolic subgroup of $G$, which is conjugate to $P$
under the real structure on $G$. Then $P\cap\overline{P}$ is a subgroup of $G$
which is defined over $\R$ and must contain a maximal torus $S$ defined over $\R$
[2]. As a maximal torus of the intersection of two parabolic subgroups of $G$ is
a maximal torus of $G$ [2], $S$ is a maximal torus in the collection
$\mathcal{T}$.

By the one-parameter subgroup $\lambda_S$ of $S$, we have an action of $\C-\{0\}$
on $V$ such that $v$ is not stable for this action and the maximal compact
subgroup of $\C-\{0\}$ preserves the Hermitian norm on $V$. As the statement
($*$) for $\C-\{0\}$ implies the statement ($*$) for $G$, it will suffice to
prove Theorem 0.2 when $G=\C-\{0\}$. Better yet, we may even assume that
$\lim_{t\to 0}t*v$ exists in $V$ and does not equal $v$, where $*$ denotes the
action of $\C-\{0\}$ on $V$.

As in section 2, we may write $p_v'$ on
$\text{maximal compact}\backslash(\C-\{0\})=\R$ uniquely
in the form $\sum a_i e^{\ell_i x}$ with positive $a_i$'s and increasing real
$\ell_i$'s. As the above limit exists, the $\lim_{x\to-\infty}p_v'(x)$ exists in
$\R$ and, hence, the $\ell$'s are not negative.
% Note: x = log|t| and t -> 0, so the relevant limit is x -> -infinity
As the above limit does not equal $v$, at least one $\ell_i$ must be positive.
Thus, $p_v'(x)$ is a strictly increasing function on $\R$. Hence, $p_v'$ and
$p_v$ never obtain a minimum value. This proves the statement ($*$) and, hence,
Theorem 0.2 is true.
\end{proof}

To finish our criteria for stability, it remains to dramatize the meaning of $v$
being properly stable in terms of the function $p_v$. We will do this in

\renewcommand{\thethm}{4.1}
\begin{thm}
If $v$ is properly stable, then
\begin{abc}
\item the induced function $p_v'$ on $K\backslash G$ is a Morse function with one
      critical point where it obtains its minimum value.
\end{abc}
If furthermore $p_v$ obtains its minimum value at the identity $e$, then
\begin{abc}[start=2]
\item the stabilizer $G_v$ is contained in $K$.
\end{abc}
Conversely, if $p_v'$ obtains a minimum at a unique point of $K\backslash G$,
then $v$ is properly stable.
\end{thm}

\begin{proof}
For a), note that Theorems 0.1 and 0.2 say that the only critical points of $p_v$
are situated on a connected $K-G_v$ coset $m$ where $p_v$ obtains its minimum
value. As $G_v$ is finite, this double coset $m$ is actually a $K$-coset. Hence,
$p_v'$ has only one critical point on $K\backslash G$ where it obtains its
minimum. The last part of Theorem 0.1 means in this case that the second order
variation of $p_v'$ in any non-zero direction through its critical point is
positive. In other words, $p_v'$ is a Morse function. This proves a).

To prove b), note that a) implies that $K\cdot G_v=K$ as $K$ is the unique
critical point of $p_v'$ and $p_v'$ is right invariant under $G_v$. Therefore,
$G_v\subseteq K$. This proves b).

For c), note that Theorem 0.2 shows that $v$ is stable because $p_v$ obtains a
minimum value. As we may assume without loss of generality that $p_v'$ obtains
its minimum at the point $K$ in $K\backslash G$, we have the inclusion
$K\cdot G_v\subseteq K$. Hence, $G_v\subseteq K$. As $G_v$ is a compact algebraic
subgroup of $G$, it can only be contained in the compact group $K$ if $G_v$ is
finite. Therefore, $G_v$ is finite and, hence, $v$ is properly stable.
\end{proof}

Trivially we want to mention a special way to find a vector in the orbit
$G\cdot v$ where $p_v$ obtains a minimum value.

\renewcommand{\thecor}{4.2}
\begin{cor}
Assume that $v$ is properly stable and $G_v$ is contained in $K$. If
$K=\{g\in G\mid gG_vg^{-1}\subseteq K\}$, then $p_v$ obtains its minimum value at
$e$.
\end{cor}

\begin{proof}
The stabilizer of $g\cdot v$ is $gG_vg^{-1}$. By Theorem 4.1, we know that
$gG_vg^{-1}\subseteq K$ if $p_v$ has a minimum at $g$. Therefore the hypothesis
means that the coset $K$ is the only possibility for the set where $p_v$ obtains
% CORRECTION: "possiblity" -> "possibility"
its minimum.
\end{proof}

An example where this hypothesis is satisfied is when $G=\mathrm{SL}(W)$, the
induced representation of the finite group $G_v$ on $W$ is irreducible and $K$
consists of the elements of $G$ which preserve some Hermitian norm of $W$.
% CORRECTION: "consists the elements" -> "consists of the elements"
The hypothesis applies because of the uniqueness up to scalar of the Hermitian
norm invariant under the action of an irreducible representation.

This kind of example actually occurs in the study of covariants of abelian
varieties over $\C$ embedded by complete linear systems.

\bigskip
\begin{flushleft}
\textbf{References}
\end{flushleft}
\begin{list}{}{\setlength{\leftmargin}{2.5em}\setlength{\labelwidth}{2em}%
  \setlength{\itemindent}{-2.5em}\setlength{\itemsep}{0.6ex}}
\item[{[1]}] D.~Birkes, Orbits of linear algebraic groups, Annals of Mathematics
  93(1971), 459--475.% CORRECTION: "algebriac" -> "algebraic"
\item[{[2]}] A.~Borel, Linear Algebraic Groups, Benjamin, New York, 1969.
\item[{[3]}] A.~Borel, Introduction aux groupes arithmétiques, Hermann, Paris,
  1969.% CORRECTION: added accent: "arithmetiques" -> "arithmétiques"
\item[{[4]}] A.~Borel and Harish-Chandra, Arithmetic subgroups of algebraic
  groups, Annals of Mathematics 75(1962), 485--535.
\item[{[5]}] S.~Helgason, Differential Geometry and Symmetric Spaces, Academic
  Press, New York, 1962.
\item[{[6]}] G.~Kempf, Instability in invariant theory, Annals of Mathematics
  108(1978), 299--316.% CORRECTION: title "Instability in invariant theory"; "to appear" -> published coordinates
\item[{[7]}] G.~Mostow and T.~Tamagawa, On the compactness of arithmetically
  defined homogeneous spaces, Annals of Mathematics 76(1962).
\item[{[8]}] D.~Mumford, Geometric Invariant Theory, Ergebnisse der Math. (34),
  Springer-Verlag, Berlin, 1965.
\end{list}

\bigskip\bigskip
\begin{flushright}
\begin{tabular}{l}
The Johns Hopkins University\\
Baltimore, Maryland 21218\\
USA\\[1.5em]
University of Washington\\
Seattle, Washington 98105\\
USA
\end{tabular}
\end{flushright}

\end{document}
